\subsection{Why the Residual Is Informative}

The interest of T1 is not that it explains how midge swarms are generated or
that it defines a universal collective-motion law. Its value is that it makes
a specific layer of motion measurable. In biological groups, raw velocity
differences across space are expected because groups turn, expand, contract,
and rearrange locally. The present result is more specific: a local tangential
residual remained measurable after ordinary local affine deformation had been
admitted into the baseline and subtracted.

This places T1 between whole-swarm state variables and individual behavioral
rules. It is not a raw trajectory feature, and it is not a complete mechanism.
It is a bounded local residual layer whose survival, form, and explanatory
limits can be tested observation by observation.
T1 should therefore not be interpreted as a force, a plastic rearrangement
variable in the materials-science sense, or a direct interaction metric. The
non-affine construction is used here only as a kinematic diagnostic of local
relative motion not represented by the fitted affine map.

The positive and negative results therefore work together. Survival after
local affine subtraction establishes a phenomenon worth explaining. The failure
of the tested compact-state, event-timing, and recent-history reductions shows
that several natural simple explanations were insufficient under their current
definitions. The study thus identifies a measurable empirical target and the
boundaries of several candidate reductions.
